\section{Discussion}
\label{sec:discussion}

\subsection{What distinguishes \ours{} from related ensemble methods?}

The closest existing methods are snapshot ensembling (cyclic LR within one
run) and SWA (weight averaging along one trajectory). Our method differs
from both:

\textbf{vs.\ Snapshot ensembling.} Snapshot ensembling runs \emph{one}
training with cyclic LR restarts, collecting snapshots at each local
minimum. Our method runs \emph{two separate} trainings with different
$T_\mathrm{max}$, no cycles. Conceptually simpler (no LR modification) and
preserves the original training recipe.

\textbf{vs.\ SWA.} SWA averages \emph{weights} collected along a single
cosine trajectory after the LR has decayed. Our method averages
\emph{outputs} across two trajectories. Output averaging preserves the
ensemble's calibration and uncertainty properties (weight averaging
collapses them into a single point estimate). Output ensembles also
handle multiple minima better: our mode-connectivity probe
(\cref{sec:basin}) shows that the linear weight-space path between a
100ep checkpoint and a 150ep checkpoint stays at chance for every
interpolation coefficient $\lambda \in [0, 1]$. In other words,
parameter-space averaging between the two groups simply does not work,
while prediction-space averaging yields a consistent $+1.5$~pp
improvement --- a textbook separation between the two strategies.

\textbf{vs.\ Fast Geometric Ensembles.} FGE explicitly walks curves
between minima in weight space, requiring a pretrained starting point and
a curve-finding loss. Our method achieves comparable diversity by simply
varying the endpoint of an already-used schedule---no extra training
infrastructure, no new loss.

\textbf{vs.\ Standard deep ensembles.} Deep ensembles differ only in seed.
Our method adds \emph{trajectory length} as a second diversity axis on
top of seed. The $+1.5$~pp gain we observe is on top of what a standard
5-seed deep ensemble achieves.

\subsection{Why is the gain largest for capacity-bound students?}

Our experiments show a clean scaling: gain is largest on hard tasks with
low individual accuracy. The mechanism is simple. At low individual
accuracy, many test samples are near the decision boundary---the model
is uncertain and its prediction is determined by small features. Different
training trajectories find slightly different decision boundaries, and
ensembling corrects errors on the samples where they disagree.

On easy tasks (MNIST at 0.985 individual accuracy), essentially every
test sample is confidently correct, so there are few disagreement samples
and little gain. At the extreme of trivial tasks, ensembling gives zero
gain (everything is unanimous and correct).

\subsection{When should practitioners use \ours?}

We recommend \ours{} when:
\begin{itemize}[leftmargin=*]
\item \textbf{Individual accuracy is in the 0.5--0.85 range.} This is where
the gain is largest.
\item \textbf{You already use a cosine LR schedule.} No modifications
needed, just vary $T_\mathrm{max}$.
\item \textbf{Compute budget permits a second training run.} The trick
doubles training compute, though you can use smaller $T_B - T_A$ to
reduce this.
\item \textbf{Calibrated ensemble outputs matter.} Unlike SWA, \ours{}
produces a true probabilistic ensemble, preserving uncertainty and
calibration.
\end{itemize}

We \emph{do not} recommend it when individual accuracy is near ceiling
(MNIST-level tasks) or when compute is so tight that a second training
run is infeasible.

\subsection{Limitations}

\paragraph{Requires a second training run.} \ours{} costs roughly
$(T_A + T_B)/T_A - 1$ extra compute on top of a standard
deep ensemble. For $T_A=100, T_B=150$ this is $0.5\times$ extra compute.
Not free.

\paragraph{Endpoint ratio has an operating regime.} Our endpoint sweep
(\cref{sec:endpoint_sweep}) shows that $\rho = T_B/T_A \approx 1.5$ is
a sweet spot: smaller ratios ($\rho \in [1.1, 1.4]$) yield smaller gains,
and much larger ratios ($\rho \approx 2.0$) cause group $B$ to overfit,
pulling the mixed ensemble below the $T_A$-only baseline. We recommend
$\rho \in [1.4, 1.6]$; outside this window the user should expect
reduced or no gain.

\paragraph{Same-seed, different-length vs.\ different-seed, different-length.}
We use different random seeds for each (seed, length) pair. We have not
tested whether the trick also works with the \emph{same} random seed for
the 100ep and 150ep runs (only LR schedule differs). The non-determinism
in SGD data ordering suggests this would still give some diversity.

\paragraph{Ceiling-bound datasets give smaller gains.} The scaling law
is also a scope statement: on tasks where individual accuracy is near
ceiling (MNIST at 0.985, SST-2 at 0.91) the headroom is small and the
gain is correspondingly small. We view this as predictive rather than
limiting---the scaling fit tells practitioners in advance when the
trick will and will not pay off.
